\section{Cuando el dominio no permite verificación rápida: entornos lentos, reward model y descomposición}

\videofigure{figures/self-improvement-directions.jpg}{La autosuperación futura necesita una generalización entre dominios más sólida, verificación robusta/meta, así como una selección automática de los datos de train-time scaling.}{00:33:34--00:40:08}

Las matemáticas y el código tuvieron éxito primero porque el reward es rápido e inequívoco. El diseño de chips con simulación, los experimentos científicos, la robótica y las tareas creativas pueden requerir horas, días o un juicio subjetivo, por lo que no se pueden ejecutar decenas de miles de veces dentro del RL.

\begin{importantbox}{No poder verificarse rápido no equivale a no poder verificarse en absoluto}
Se puede convertir una evaluación real costosa en una señal gold escasa y luego entrenar un reward model surrogate; también se puede descomponer una tarea grande en subverificadores más rápidos como compilación, simulación local o perfilado de desempeño. Pero la señal proxy debe calibrarse periódicamente con el entorno real.
\end{importantbox}

Una forma de aproximación consiste en aprender un reward model $\hat R_\phi(x,y)$ que prediga el retorno del entorno costoso $R(x,y)$:

$$
\phi^*=\arg\min_\phi
\mathbb E_{(x,y,R)\sim \mathcal D}
\left[\hat R_\phi(x,y)-R(x,y)\right]^2.
$$

Después, el RL itera rápidamente sobre $\hat R_\phi$, pero la política buscará activamente los puntos ciegos del reward model, por lo que se necesita active sampling para devolver las muestras de alta incertidumbre a la simulación real o a expertos.

\begin{warningbox}{El alcance de generalización del surrogate reward determina el límite superior de seguridad}
Si la política se aleja de la distribución de los datos de entrenamiento, el reward model puede otorgar una puntuación alta a diseños físicamente inválidos o peligrosos. Cuanto mejor sabe optimizar el modelo, más amplifica los pequeños errores del evaluador.
\end{warningbox}

\begin{knowledgebox}{La descomposición de tareas es un puente de capacidad, no la respuesta final}
Cuando el modelo actual no puede optimizar un chip completo, se puede empezar con el kernel, los módulos locales y el análisis de perfilado. La descomposición permite que las capacidades existentes participen en flujos complejos, pero al final aún hay que verificar si la combinación de las subtareas satisface el objetivo global.
\end{knowledgebox}

\subsection{Resumen de la sección}

En dominios de verificación lenta y subjetiva, la autosuperación necesita un reward surrogate, tareas jerarquizadas y retroalimentación real escasa. El reto clave pasa de «¿hay reward?» a «¿el reward proxy sigue siendo confiable después de que la política se optimizó con él?».
